\documentclass[11pt,a4paper]{article}
\usepackage[T1]{fontenc}
\usepackage{lmodern,amsmath,amssymb,amsthm,microtype}
\usepackage[margin=27mm]{geometry}
\usepackage[colorlinks=true,linkcolor=blue,urlcolor=blue,citecolor=blue]{hyperref}
\setlength{\emergencystretch}{2em}
\newtheorem{theorem}{Theorem}[section]
\newtheorem{proposition}[theorem]{Proposition}
\newcommand{\Q}{\mathbb Q}
\newcommand{\Z}{\mathbb Z}
\title{A quadratic obstruction to retaining half a scalar helper}
\author{}\date{17 November 2026}
\begin{document}\maketitle
\begin{abstract}
For a specified translated-secant chart, retaining only the rational
component of the product of the scalar quotient and the secant intercept
does not give a quadratic presentation of the scalar-quotient graph.
An exact finite-field minor certifies the complete space of its quadratic
identities in characteristic zero. All these identities admit an explicit
wrong nonzero scalar quotient at a nontorsion point pair. The result
concerns this interface and this helper, not the full 195-coordinate
system or arbitrary alternative charts.
\end{abstract}

\section{The interface and its guards}

Put $K=\Q(\rho)$, $\rho^2=2$, $D=11-6\rho$, and
$E:y^2=x^3-Dx$. Use six field coordinates $F,G,s,t,h,d$, and
write $p=h+s+t$. The translated-secant chart is
\begin{equation}\label{eq:chart}
 d=pF,\qquad d=hG,\qquad
 D(F+G)=(s-t)h-t^2,\qquad
 1=d(s-t)-DFG.
\end{equation}
Adjoin the already available square outputs $F_2=F^2$, $G_2=G^2$.
For the present argument restrict to the nonempty open set
\[
 FG(F-G)stdS\ne0,\qquad d(h+s)-1\ne0,
\]
where $S\in K$ is a free nonzero reference coordinate. A field
nonzero condition means that its rational norm is nonzero.
The chart recovers
\[
 U=(F^{-1},s/F),\quad V=(G^{-1},t/G),\quad
 x(U-V)=hp,\quad q=-\frac{y(U-V)}{x(U-V)}=h+s-d^{-1}.
\]
These formulas also follow directly from the secant through the two
translated points. On the displayed open set all denominators used
below are nonzero.

The scalar-quotient graph requires
\begin{equation}\label{eq:scalar}
                  Z=-2S/q,
 \qquad Z\bigl[d(h+s)-1\bigr]+2Sd=0.
\end{equation}
Thus $Z\ne0$. One quadratic presentation uses a field helper $e=Zd$
and the field equation $e(h+s)-Z+2Sd=0$. Consider instead retaining
only $e_0$, the rational component of $e=e_0+e_1\rho$.
There are then 21 rational coordinates: the two components of
\[
              F,G,s,t,h,d,S,Z,F_2,G_2
\]
and $e_0$. No other reference, ratio, normalizer or scalar-square
coordinates of the 195-coordinate source are allowed in this interface.

The eliminated component has the quadratic expression
$e_1=Z_0d_1+Z_1d_0$. Its two consumers become
\begin{align*}
 e_0(h_0+s_0)+2(Z_0d_1+Z_1d_0)(h_1+s_1)
       -Z_0+2(S_0d_0+2S_1d_1)&=0,\\
 e_0(h_1+s_1)+(Z_0d_1+Z_1d_0)(h_0+s_0)
       -Z_1+2(S_0d_1+S_1d_0)&=0.
\end{align*}
These are cubic expressions. The following theorem shows that replacing
them by other quadratic identities in the same interface cannot repair
the loss.

\section{The complete quadratic identity space}

\begin{theorem}\label{thm:space}
The vector space over $\Q$ of polynomials of degree at most two
vanishing identically on the scalar-quotient graph has dimension 17.
A basis consists of the two rational components of the four equations
in \eqref{eq:chart}, of
\[
 F_2-F^2=0,\quad G_2-G^2=0,\quad
 Fd-pF_2=0,\quad Gd-hG_2=0,
\]
and the one rational equation
\[
                     e_0-Z_0d_0-2Z_1d_1=0.
\]
In particular, no quadratic identity imposes the scalar quotient after
these chart identities and the definition of $e_0$ have been imposed.
\end{theorem}
\begin{proof}
The displayed 17 equations hold identically. The last two field
equations follow by multiplying $d=pF$, $d=hG$ by $F,G$ and using
the square outputs. Their coefficient vectors have rank 17 modulo
1009, as checked by the retained exact certificate, so they are
linearly independent over $\Q$.

There are $\binom{23}{2}=253$ monomials of degree at most two
in 21 coordinates. The certificate evaluates them at 300 points
of the guarded scalar graph modulo 1009, using $\rho=439$.
The evaluation matrix has a $236$ by $236$ minor with determinant
\[
                             514\pmod {1009}.
\]
The selected rows, columns and evaluation points are retained. The
17 displayed relations annihilate every row. Thus the evaluation
matrix has rank exactly 236 and kernel dimension 17.

Here reduction modulo 1009 bounds characteristic-zero identities,
not merely identities at a finite selection of unrelated points.
To see this, parameterize a dense open part of the graph by the
two elliptic points $U,V$ and the free reference $S$ over $K$.
After extending scalars to $K$, restriction of scalars gives four
independent elliptic-point slots, two on $E$ and two on its conjugate,
and two free reference slots. Work over
$A=\Z_{(1009)}[\rho]$ and invert 2, $\rho$, and the denominators
of these rational coordinate functions. The affine equations for
the four elliptic slots are monic in their separate squared
ordinates. Their coordinate ring, with the free reference variables,
is therefore torsion-free over $\Z_{(1009)}$; localization preserves
this property. All selected points have the inverted denominators
nonzero, and the curves have good reduction at 1009.

Consequently every characteristic-zero polynomial identity with
coefficients in $\Z_{(1009)}$ vanishes in this localized ring and
specializes to zero at every selected point. If $W\subset\Q^{253}$
is the space of characteristic-zero identities, then
$W\cap\Z_{(1009)}^{253}$ is saturated. Since
$\Z_{(1009)}$ is a discrete valuation ring, its reduction has
dimension $\dim_\Q W$ and lies in the 17-dimensional evaluation
kernel. Therefore $\dim_\Q W\le17$, proving the theorem.
\end{proof}

\section{An admissible wrong quotient}

\begin{proposition}
The quadrics of Theorem~\ref{thm:space} do not define the scalar
graph on its stated nonzero domain, even at a nontorsion point pair.
\end{proposition}
\begin{proof}
Take the retained point $R_0=(1+\rho,2\rho)$ and $P=2R_0$,
where
\[
 P=\left(\frac{57}{8}-\frac72\rho,
           -\frac{165}{8}+\frac{417}{32}\rho\right).
\]
Set $U=R_0$, $V=P$. Their difference is $-R_0$, which is
nontorsion under the retained fixed-curve group assertion. The
chart coordinates are
\begin{align*}
 F&=-1+\rho,&G&=\frac{456+224\rho}{1681},\\
 s&=4-2\rho,&t&=-\frac{87}{41}+\frac{129}{164}\rho,\\
 h&=-1+\frac74\rho,&d&=\frac{8+14\rho}{41}.
\end{align*}
Here $hp=1+\rho$ and $q=s$. Choose $S=s$.
The required scalar quotient is $Z=-2$. Instead set $Z=1$,
$e_0=8/41$, $F_2=F^2$, and $G_2=G^2$.
Every listed nonzero guard holds, and every one of the 17 quadratic
identities vanishes. However the residual in \eqref{eq:scalar} is
\[
                       \frac{-72+120\rho}{41}\ne0.
\]
Thus this point satisfies every quadratic identity of the graph and
does not belong to the graph.
\end{proof}

The rationally simple restriction of scalars and the nontorsion
point $R_0$ imply that $E(K)$ is Zariski dense in
$\operatorname{Res}_{K/\Q}E$, as in \cite{Families}. Products and
nonempty open subsets preserve this density. Hence requiring a
quadric to vanish at every rational point of this interface graph
already makes it an identity, so the obstruction also applies to
an exact rational zero-set presentation. Adding finitely many
nonvanishing polynomial guards cannot fix the interface: after fixing
any admissible base chart and $S$, all 17 quadrics allow arbitrary
$Z\in K$, with its prescribed $e_0$. Guards that hold at the true
quotient leave a nonempty rational open set of wrong quotients.

\section{The separate torsion boundary}

Even retaining both components of $e=Zd$ requires attention to the
scalar source's nontorsion condition. For $U=R_0$ and
$V=R_0+(0,0)$, the same six-field chart has
\[
 F=-1+\rho,\quad G=-\frac{23+17\rho}{49},\quad
 s=4-2\rho,\quad t=-s,\quad h=d=0.
\]
Its difference is $(0,0)$. The polynomial helper equations admit
$Z=e=0$ for every $S$, although this branch is excluded from the
inverse-Vieta scalar source. The nonzero guard on $d$ removes this
boundary; it does not remove the nontorsion wrong-quotient example.

The exact checks use only the standard Python library. They verify
the rational examples, all finite-field graph equations, the 17
coefficient relations, the full evaluation rank and the displayed
nonzero minor. They neither assert a lower bound for the complete
195-coordinate construction nor rule out use of its additional global
coordinates, a different helper, or a different polynomial presentation.

\begin{thebibliography}{9}
\bibitem{Families}\emph{Curves and abelian families in an actual-trace square locus},
\nolinkurl{negative_results/papers/pole_parameters/curves_and_abelian_families_20261117_v69.tex}.
\bibitem{Scalar}\emph{A quartic in 198 variables},
\nolinkurl{sources/h10_quartic198/arithmetic/inverse_vieta_198.tex}.
\end{thebibliography}
\end{document}
